\section{Related work}
\label{sec:related}

\paragraph{Language models as agents and decision-makers.}
SmartPlay \citep{wu2024smartplay} probes nine capabilities with six games, one of them a two-armed bandit, and finds that models differ by capability rather than by a single agentic skill. BALROG \citep{paglieri2025balrog} grades difficulty from BabyAI to NetHack; models make progress on the easy environments and fail on the hardest. LMAct \citep{ruoss2025lmact} shows that even 512 expert episodes in context rarely bring frontier models to expert level on tasks that include tic-tac-toe and grid navigation. AgentBench \citep{liu2024agentbench} evaluates 29 API and open models in eight environments and names weak long-term reasoning, decision-making and instruction following as the main obstacles, and GTBench \citep{duan2024gtbench} finds that LLMs fail in complete-information deterministic games such as tic-tac-toe. \citet{huang2022zeroshot} showed that plans generated by language models are often not admissible actions, the problem that likelihood scoring avoids by construction. These benchmarks establish that difficulty depends on task structure, the premise of the task-dependent slope in H1, but their model sets are not controlled ladders and each fixes one conversion method. The design nearest to ours is that of \citet{jiwatode2026spatial}, who cross Qwen3 model size with thinking mode and planning horizon on spatial games; we extend it to several families, five conversion methods and six decision structures.

\paragraph{Exploration and bandits.}
\citet{krishnamurthy2024explore} found that, across many prompt designs, only GPT-4 with chain-of-thought and an externally summarized history explored a multi-armed bandit reliably. EVOLvE \citep{nie2025evolve} relates regret to model size and shows that algorithm-guided inference and algorithm distillation let smaller models explore better than larger ones without such support. \citet{schmied2026greedy} trace the failures of 2B to 27B models to greediness, with up to 55\% of the action space left unexplored, to a frequency bias that fades with scale, and to a knowing-doing gap: rationales were correct about 87\% of the time, yet the chosen action often did not follow from them. \citet{harris2026explore} find that LLMs often fail even to exploit, falling short of simple linear regression, and \citet{monea2025icrl} show that models from 500M to 70B parameters can learn from reward in context. Algorithm distillation bears on how such policies should be trained: cloning the learning histories of a reinforcement-learning algorithm, not expert actions, is what teaches in-context reinforcement learning \citep{laskin2023incontext}. We use UCB1 \citep{auer2002ucb} as a realizable reference and standard regret definitions \citep{lattimore2020bandit}. H6, that exploration is poor without scaffolding, is therefore largely established. What we add is its dependence on scale within families, its behavior under conversion methods other than prompting, and a separation of exploration errors from exploitation errors.

\paragraph{Scale, emergence and small models.}
Pre-training loss falls predictably with parameters and training tokens \citep{hoffmann2022chinchilla}, which motivates the log-linear form of H1. The same work complicates any reading of such a slope, because the small rungs of modern ladders are trained far beyond the compute-optimal number of tokens; SmolLM2-1.7B, for example, saw about 11T tokens \citep{allal2025smollm2}. Parameter count therefore partly stands in for tokens per parameter. Emergent abilities \citep{wei2022emergent} are the main alternative to a smooth slope, and \citet{schaeffer2023mirage} attribute many of them to discontinuous metrics, one reason we fit slopes to continuous normalized returns rather than to success rates. Few-shot ability improves steadily with scale within one recipe \citep{brown2020gpt3}, and the length of task a model can execute keeps growing with size after single-step accuracy has saturated \citep{sinha2026illusion}. The case for small models rests on extra computation and adaptation. In FLOPs-matched comparisons a small model given more test-time compute can beat one 14 times larger on problems it already solves some of the time \citep{snell2025ttc}, and algorithm distillation lets small models explore bandits better than larger unaided ones \citep{nie2025evolve}, which partly answers the in-distribution half of H3 for that task family.

\paragraph{Converting a language model into a policy.}
Scoring answer strings by probability is distorted by surface-form competition between paraphrases \citep{holtzman2021surface}. Scoring an option symbol works only for models that bind symbols to options well \citep{robinson2023mcsb}, option identifiers carry a token-level prior \citep{zheng2024selectors}, and few-shot predictions lean toward majority, recent and common labels \citep{zhao2021calibrate}. For instruction-tuned models, first-token probabilities often disagree with the answer the model writes \citep{wang2024myanswerc}. Meaning-preserving changes of prompt format move few-shot accuracy by up to 76 points \citep{sclar2024formatting}, and the serialization of tabular inputs matters for LLM classifiers \citep{hegselmann2023tabllm}. As a policy, likelihood scoring of legal actions underlies GLAM \citep{carta2023glam} and TWOSOME \citep{tan2024twosome}, and the latter shows that long action strings are penalized unless scores are normalized for length. This work motivates H4, but it was measured on question answering, not on sequential decisions across a size ladder.

A linear probe on frozen hidden states shows what is decodable, which is not the same as what the model uses, and needs control tasks to be interpreted \citep{belinkov2022probing}. Intermediate layers are often better features than the last \citep{skean2025layer}, and models often encode the correct answer while generating a wrong one \citep{orgad2025know}. That gap is how a probe could beat generation on our tasks, and it is why we pair the probe with a control trained on raw task features.

Behavior cloning \citep{pomerleau1988alvinn} compounds its errors once the learner leaves the expert's state distribution \citep{ross2011dagger}, and LoRA \citep{hu2022lora} and QLoRA \citep{dettmers2023qlora} make it affordable at every rung of a ladder. Whether a cloned policy generalizes is contested. Supervised fine-tuning memorizes the rules it was trained on while outcome-reward RL generalizes \citep{chu2025sft}, which supports the out-of-distribution half of H3; frozen language models with RL-trained adapters \citep{szot2024llarp} and LoRA fine-tuning on human choice data \citep{binz2025centaur} generalize to some shifts, which cuts against it. LIMA \citep{zhou2023lima} and URIAL \citep{lin2024urial} suggest that alignment mostly changes style and format. That predicts the second clause of H2, that base and instruct checkpoints converge under probing and fine-tuning, but it also threatens the first, since a base model with a good few-shot prompt may close the gap under prompting as well. DPO \citep{rafailov2023dpo}, GRPO \citep{shao2024deepseekmath} and self-play on text games \citep{liu2026spiral} are conversions that \deadeye{} v1 leaves out.

\paragraph{Reasoning and output format.}
Chain-of-thought helps large models, while small ones produce fluent but illogical chains \citep{wei2022cot}. More reasoning can also hurt. Deliberation lowers accuracy on tasks where it hurts people too, such as implicit statistical learning \citep{liu2025mindstep}; on controllable puzzles, thinking helps only at intermediate complexity \citep{shojaee2025illusion}; longer reasoning can reduce accuracy \citep{gema2025inverse}; and models of about 3B parameters or fewer do not reliably benefit from long chains \citep{li2025smallstruggle}. Strict output formats can suppress reasoning \citep{tam2024speakfreely}, so our chain-of-thought prompt lets the model reason freely before a delimited action line. Every clause of H7 finds support somewhere in this work, but on different tasks and models. Switching thinking on and off inside the same weights \citep{yang2025qwen3} permits a cleaner test, and latency-normalized efficiency has not been compared across sizes and task structures.

\paragraph{Quantization.}
At a fixed total bit budget, 4-bit weights are close to optimal on zero-shot language tasks \citep{dettmers2023kbit}, and GPTQ \citep{frantar2023gptq} and AWQ \citep{lin2024awq} quantize large models with little loss. Precision-aware scaling laws predict that the penalty for post-training quantization grows with the amount of pre-training data \citep{kumar2025precision}; small modern models are the most over-trained, so this is the mechanism behind H5. These studies measure loss, perplexity or accuracy on language benchmarks, so for H5 there is a mechanism but no direct evidence on multi-step decisions.

\paragraph{Evaluation methodology.}
We follow \citet{agarwal2021precipice} in reporting interquartile means and bootstrap confidence intervals, and \citet{henderson2018matters} in fixing seeds and hyperparameter budgets across all rungs. Because results for language models depend heavily on prompt format and on the choice between scoring and generation \citep{sclar2024formatting,wang2024myanswerc}, we report both conversions side by side and release every decision.

\paragraph{Cognitive and behavioral studies of decisions.}
Studies in the style of experimental psychology ask how models decide, not only how well. GPT-3 outperformed people on a two-armed bandit but showed no directed exploration \citep{binz2023cognitive}. Even GPT-4 falls short of human-like rationality when it must form and keep beliefs \citep{fan2024rational}, language models show human-like risk and loss aversion \citep{jia2024decision}, and they assume that people are more rational than people are \citep{liu2025rational}. Our oracles are normative rather than human, but these findings suggest examining the direction of errors, for instance whether a model stands too early in blackjack, and not only their rate (Section~\ref{sec:analysis}).

\paragraph{Decision models.}
While this study was being prepared, ``decision models'' became a product category. TypeSafe AI's Jev, announced in September 2026, takes a state and typed questions and returns a chosen option with a probability for every allowed answer and a confidence, never prose; the company describes it as a ``System One'' model trained on synthetic data with a method it calls reinforcement learning for calibrated decisions, and keeps the weights and the architecture closed \citep{typesafe2026jev,typesafe2026systemone,typesafe2026choice,techcrunch2026jev}. Open-weight relatives followed within weeks. Kev attaches a rank-16 LoRA adapter and a small pointer head to Qwen3.5 and Qwen3.8 bases at 0.8B to 27B parameters, ships under Apache-2.0 with a server that speaks the same protocol, and reports accuracy close to Jev's on held-out classification sets while noting that Jev's training data is unknown \citep{palmer2026kev}. Cloudflare's Clef and Clef-flash freeze Qwen3.8-27B and Qwen3.5-9B and train a routing head with low-rank adapters, are served on Workers AI with vendor-reported median latencies of about 209 and 39 milliseconds against about 524 for Jev, and are likewise Jev-compatible \citep{cloudflare2026clef}; Perplexity's pplx-decider fine-tunes Qwen3.8-27B and reports a small lead over Jev on an eleven-task panel \citep{perplexity2026decider}; Amazon released a 2B decider of its own \citep{techcrunch2026decisionmodels,thenewstack2026jev}. In the terms of this paper these systems are the probe and LoRA conversions productized: an open base, frozen or lightly adapted, plus a head that reads probabilities over the allowed answers, exposed through an interface that is exactly our likelihood scoring. What they have not been evaluated on is decision quality against a known optimal policy; their published numbers are label accuracies on classification panels, mostly measured by the vendors themselves. \deadeye{} treats them as one more backend and places them on the same tasks, seeds and normalized scale as the conversions it studies (Section~\ref{sec:setup:decision}).

